% ============================================================
%  INTRODUCTION: Meera's dollar problem        (Member 1)
%  Main-text words: see word count in README (boxes excluded)
% ============================================================
\section*{Introduction: Meera's dollar problem}
\addcontentsline{toc}{section}{Introduction}

In October 2026, Meera, a Class 12 student in Chennai, got the email her family had been waiting for: a college in the United States had offered her a place. The first-year fees were US\$40,000, exactly what the college had charged a year earlier. Yet when her father asked his bank what the fees would cost in rupees, the answer was about ₹2.5 lakh more than it would have been a year before.\sidenote{The fees are hypothetical. The exchange rates are real closing rates: ₹88.67 per US dollar on 13 October 2025 (Business Standard, 2025) and ₹95.05 on 15 October 2026 (HDFC Sky, 2026). 40,000 × (95.05 − 88.67) = ₹2,55,200.} The college had not raised its price. The price of the dollar had gone up.

Meera's uncle Ravi, an electrician in Dubai, had the opposite surprise. Every month he sends part of his salary home to his parents in Thanjavur. This year the same amount arrived as more rupees than before.

One family, one price, and two very different outcomes. This chapter explains what that price is, who sets it, and why it keeps moving.

\begin{keypoints}{Questions this chapter answers}
\begin{itemize}
  \item Why does one country's money have a price in another country's money?
  \item Who buys and sells currencies, and why?
  \item What decides that price, and why does it change so often?
  \item How do countries, and the Reserve Bank of India, manage their exchange rates?
  \item Who gains and who loses when the rupee rises or falls?
\end{itemize}
\end{keypoints}
